% TOP_PROBLEM: {"id": 2302056, "title": "Research Problems in Function Theory — Problem 2.56", "category_id": 8, "category": {"id": 8, "name": "analysis", "display_name": "Analysis", "description": "Limits, continuity, calculus, and function theory.", "slug": "analysis", "order_index": 8, "created_at": "2026-07-31T15:26:25.671Z"}, "status": "open"}
% FIRST_POSTED: null
% ENTRY_KIND: statement_only
% Imported from: batch11.tex; UnsolvedMath ID 2302056
% Compile twice: lualatex 2302056.tex
\documentclass[11pt]{article}
\usepackage{amsmath,amssymb,mathtools}
\usepackage{fontspec}
\usepackage{unicode-math}
\setmainfont{Latin Modern Roman}
\setmathfont{Latin Modern Math}
\usepackage{xurl}
\usepackage[unicode,colorlinks=true,linkcolor=blue,urlcolor=blue,pdftitle={UnsolvedMath Problem 2302056}]{hyperref}
\newcommand{\sref}[2]{\hyperlink{source:#1:#2}{[S#2]}}
\newcommand{\cataloguescope}[1]{\paragraph{Mathematical statement.}#1}
\begin{document}
% UnsolvedMath ID 2302056; source catalogue code AMR-022-2056
\setcounter{section}{2302055}
\section{Research Problems in Function Theory — Problem 2.56}\label{2302056}
{\small\sffamily UnsolvedMath ID 2302056 · Analysis}
\subsection{Definitions and mathematical statement}

\paragraph{Shared notation.}$\mathbb N=\{1,2,\ldots\}$, and $\mathbb Z,\mathbb Q,\mathbb R,\mathbb C$ denote the integers, rationals, reals and complex numbers. Any other conventions are specified locally.


For an entire function $u:\mathbb C^n\to\mathbb C$, set $G_u=u(\mathbb C^n)$. Given a sequence $z_j\in\mathbb C^n$ with $\|z_j\|\to\infty$, say that $u$ has the $L$-property on this sequence if for every compact $K\subset G_u$, only finitely many $u(z_j)$ lie in $K$.
An ordered $L$-pair $(u,v)$ means that, on every such sequence, the $L$-property of $u$ implies that of $v$.
A nonconstant $f$ is an $L$-atom if every entire $u$ for which $(u,f)$ is an ordered $L$-pair can be written $u=\phi\circ f$ for a holomorphic function $\phi$ on $G_f$.
These definitions are the $\mathbb C^n$ version of the range-escape definitions in Problem 5.55.

\paragraph{Statement recorded in the supplied source.}
For $n\ge1$, is $f:\mathbb C^n\to\mathbb C$ an $L$-atom if and only if there are entire functions $f_2,\ldots,f_n$ such that
\[F=(f,f_2,\ldots,f_n):\mathbb C^n\longrightarrow\mathbb C^n\]
is a biholomorphic automorphism? \sref{2302056}{2}
The source attributes the conjecture to Rubel and states that he can prove the one-variable case. The ordering direction is the one in Problem 5.55: escape under $u$ must imply escape under the proposed atom $f$. \sref{2302056}{2}
\subsection{Short English statement}
Are the entire functions that are atomic for the range-escape relation exactly the coordinate functions of holomorphic automorphisms of complex affine space?
\subsection{Sources}
\begin{description}
\item[\hypertarget{source:2302056:1}{S1}] ULAM.AI and the UnsolvedMath Contributors, \emph{UnsolvedMath: A Curated Collection of Open Mathematics Problems}, record 2302056 (AMR-022-2056). CC BY 4.0. \url{https://huggingface.co/datasets/ulamai/UnsolvedMath}.
\item[\hypertarget{source:2302056:2}{S2}] W. K. Hayman and E. F. Lingham, \emph{Research Problems in Function Theory} (2018), Chapter 2, Problem 2.56 and Update 2.56; Problem 5.55 (L-property and L-atom). \url{https://arxiv.org/abs/1809.07200}.
\end{description}
\label{end:2302056}
\end{document}
